\newpage
\section{\textsc{Implementation}}
\label{chap:impl}
\rhead{\itshape Chapter 3: Implementation}
\lhead{}
\cfoot{\thepage}
\pagestyle{fancy}
\renewcommand{\headrulewidth}{0.4pt}

\begin{flushright}
\textit{``First, solve the problem. Then, write the code.''\\
John Johnson}
\end{flushright}

\subsection{Experimental Environment}
\label{subsec:hardware}

All local experiments ran on an Intel Core~i7-10750H (6 cores, 2.60~GHz),
16~GB DDR4 RAM, 512~GB NVMe SSD, Ubuntu~22.04 (WSL2).
Stack: Python~3.12, Docling~2.x, PyMuPDF~1.24, FAISS-CPU~1.8, instructor~1.3,
FastAPI~0.111, React~18 + Vite~5.
LLM calls (DeepSeek-V4-Flash, Gemini-2.0-Flash, Claude Sonnet) and embedding
calls (AraBERT via HuggingFace API, MiniLM-L12 locally) were made over HTTPS;
their latency is excluded from pipeline timings.

\subsection{Dataset}
\label{subsec:dataset}

The raw corpus consists of 137 historical examination documents from the Computer
Science Department of the University of Mohamed Khider Biskra, spanning five subjects:
Algorithms and Data Structures (algo), Software Engineering / Operating Systems (se),
Compilation (compilation), Statistics and Informatics for Commerce (commerce), and
Law (Law). Each subject folder contains \texttt{courses/}, \texttt{exams/}, and
\texttt{td/} sub-directories. Documents are in French, Arabic, or both.
Table~\ref{tab:dataset} reports post-processing statistics.

\begin{table}[H]
\centering
\caption{Dataset statistics after \texttt{docParser} processing.}
\label{tab:dataset}
\begin{tabular}{lrrrrr}
\toprule
\textbf{Subject} & \textbf{Raw files} & \textbf{JSON files} & \textbf{Exercises}
  & \textbf{Questions} & \textbf{Language} \\
\midrule
algo        & 36  & 25  & 71  & 195 & French \\
se          & 30  & 22  & 58  & 175 & French \\
compilation & 28  & 20  & 62  & 210 & French \\
commerce    & 14  &  8  & 24  &  58 & Arabic \\
Law         & 29  & 11  & 26  &  66 & Arabic \\
\midrule
\textbf{Total} & \textbf{137} & \textbf{86} & \textbf{241} & \textbf{704}
  & Fr.~/ Ar. \\
\bottomrule
\end{tabular}
\end{table}

\subsection{The \texttt{docParser} Pipeline}
\label{subsec:docparser}

\texttt{docParser} implements a five-stage IDP pipeline converting raw PDFs into
structured, validated, Bloom-labelled JSON records.

\subsubsection{Stage 1 --- PDF to Markdown Conversion}

Docling \cite{docling2024} converts PDFs to Markdown using TableFormer-based table
extraction and RTL-safe export. A \textit{hybrid fill} step uses PyMuPDF
\cite{pymupdf2023} to recover two failure modes: (i)~formula placeholders
(\texttt{<!-{}-~formula-not-decoded~-{}-{}>}) are replaced by 40-character anchor
search in raw text; (ii)~corrupted pseudocode/assembly inside fenced blocks is
substituted with PyMuPDF raw text. A sanitisation pass removes CJK codepoints,
null bytes, and HTML-like tags.

\subsubsection{Stage 2 --- LLM Extraction with Self-Healing Schema}

The Markdown is submitted to an LLM via \texttt{instructor} \cite{instructor2023},
which intercepts Pydantic validation errors, converts them to correction prompts,
and re-prompts up to \texttt{max\_retries=3}.
Schema: $\texttt{Exam} \supset \texttt{Exercise} \supset \texttt{Question} \supset
\{\texttt{Solution},\ \texttt{TableData},\ \texttt{DiagramData}\}$.
Question type constrained to \texttt{QCM | FREE | CODE}.
Provider priority: DeepSeek $\to$ Gemini $\to$ OpenRouter $\to$ Claude.

\subsubsection{Stages 3--5}

\textbf{Stage 3 (Grounding Validation).}
See Algorithm~\ref{alg:ground} for the full procedure.
Confidence $c$ is computed as in Equation~\ref{eq:confidence};
exams with $c < 0.70$ are flagged \texttt{needs\_review=True}.
\begin{equation}
  c = \max\!\left(0,\; 1 - 0.15|\mathcal{W}|
      - 0.5\cdot\mathbf{1}[\mathcal{E}{=}\varnothing]
      - 0.1\!\sum_{e}\mathbf{1}[\mathcal{Q}_e{=}\varnothing]\right)
  \label{eq:confidence}
\end{equation}

\textbf{Stage 4 (Bloom Classification).}
Each question is independently classified into
\textit{Factual / Conceptual / Procedural / Metacognitive}
by \texttt{gpt-oss-120b:free} at temperature~0.
An independent classifier is used rather than trusting the generator's
self-reported level \cite{meissner2024itemforge}.

\textbf{Stage 5 (VLM Diagram Transcription).}
Vector drawing clusters classified as diagrams (non-grid spatial distribution
or paths $>2$ items) are rendered at $2\times$ resolution and submitted to
Qwen2.5-VL-72B \cite{qwen2025vl} (primary) or Gemini-2.0-Flash \cite{google2024gemini}
(fallback), then linked to the nearest preceding question.

\subsection{The \texttt{exam-forge} Generation Engine}
\label{subsec:examforge}

\subsubsection{Hybrid Retrieval-Augmented Generation}

The full retrieval procedure is described in Algorithm~\ref{alg:rag}.
\textbf{Chunking:} PDFs are segmented adaptively (title-based when $\geq 2$ headings
detected, sentence-based otherwise; max 1200 chars, min 80 chars).
\textbf{Sparse:} BM25 \cite{robertson2009bm25} for exact technical term matching.
\textbf{Dense:} AraBERT \cite{antoun2020arabert} (Arabic ratio $> 0.15$) or
multilingual MiniLM \cite{reimers2019sentence} otherwise, stored in FAISS
\cite{faiss2019}.
\textbf{Fusion:}
$\mathrm{RRF}(d) = \sum_{r} \frac{1}{60 + \mathrm{rank}_r(d)}$,
capped at 2 chunks per source file.
\textbf{Cloud:} Supabase \cite{supabase2023} + pgvector \cite{pgvector2024},
scoped by \texttt{user\_id} and \texttt{subject}.

\subsubsection{Question Generation and Post-Generation Evaluation}

Three question types (MCQ, SAQ, Structured) are generated by
\texttt{deepseek/deepseek-v4-flash} \cite{deepseek2024} at temperature~0.75
with schema enforcement via \texttt{instructor}.
Style anchors (3--6 few-shot examples filtered by subject + Bloom level) and RAG
factual context are injected as separate labeled prompt blocks, preventing topic
blending (see Listing~\ref{lst:mcq_prompt}).
MCQ distractors represent common student misconceptions; SAQ includes a
\texttt{grading\_rubric} of 2--6 key-concept bullets \cite{morris2024aqg}.

Each generated question is then assessed by \texttt{openai/gpt-oss-120b:free}
on four axes via a fresh per-question RAG retrieval (anti-circularity
\cite{morris2024aqg}): \textbf{(1)} \texttt{factual\_ok} (answer derivable from
course material); \textbf{(2)} \texttt{level\_correct} (Bloom level matches
difficulty, with \texttt{corrected\_level} when wrong
\cite{meissner2024itemforge}); \textbf{(3)} \texttt{quality}
(\textit{good / needs\_revision / reject} \cite{mucciaccia2025mcq});
\textbf{(4)} \texttt{points\_calibration} and \texttt{estimated\_minutes}
for psychometric calibration \cite{isley2025quality}.
Both prompt structures are shown in Figure~\ref{lst:mcq_prompt}.

\subsubsection{Exam Planning and Assembly}

The \texttt{ExamPlanner} queries a \texttt{TemplateAnalyzer} that mines the teacher's
archive to discover recurring structural patterns (exercise type distribution,
modal point total, typical duration).
The resulting \texttt{ExamPlan} is executed by the \texttt{ExamAssembler}, streaming
each exercise to the client as an SSE event immediately upon completion.

\subsection{Web Application}
\label{subsec:webapp}

\subsubsection{Backend}

The school management backend uses FastAPI \cite{fastapi2019} and SQLite with a
modular architecture: one package per domain entity (teacher, subject, course, exam,
archive, series).
Authentication uses JWT with bcrypt hashing and Google OAuth 2.0.
A per-teacher \texttt{exam\_limit} gates AI usage.
PDF uploads trigger \texttt{parse\_pdf\_async} from \texttt{docparser\_service.py}
as a background task.

\subsubsection{Frontend}

The React+Vite interface renders questions card-by-card as SSE events arrive,
displaying each card alongside its four-axis evaluation result.
Instructors accept, edit, or reject each item before archiving the exam.

% ======================================================================
\subsection{Key Algorithms and Prompt Design}
\label{subsec:prompts}
% ======================================================================

\noindent
\begin{minipage}[t]{0.48\textwidth}
\begin{algorithm}[H]
\small
\caption{HybridRetrieve$(q,\textit{subj},k{=}10)$}
\label{alg:rag}
\begin{algorithmic}[1]
\State $r \leftarrow \text{arabic\_ratio}(q)$
\State $R_{\text{bm25}} \leftarrow \text{BM25}[\textit{subj}].\text{top}(q,k)$
\If{$r > 0.15$}
  \State $R_{\text{dense}} \leftarrow \text{FAISS\_Ar}.\text{top}(\text{AraBERT}(q),k)$
\Else
  \State $R_{\text{dense}} \leftarrow \text{FAISS\_ML}.\text{top}(\text{MiniLM}(q),k)$
\EndIf
\State $R \leftarrow \text{RRF}(R_{\text{bm25}},R_{\text{dense}},60)$
\State cap at 2 chunks/source; return top-$k$
\end{algorithmic}
\end{algorithm}
\end{minipage}
\hfill
\begin{minipage}[t]{0.48\textwidth}
\begin{algorithm}[H]
\small
\caption{GroundingValidate$(\textit{exam},\textit{md})$}
\label{alg:ground}
\begin{algorithmic}[1]
\State $W \leftarrow [\,]$
\For{$e \in \textit{exam.exercises}$}
  \If{$e.Q = \varnothing$} $W \mathrel{+}= \text{"empty"}$ \EndIf
  \For{$q \in e.Q$}
    \If{$q.\text{text}[q_1{:}q_1{+}20] \notin \textit{md}$}
      \State $W \mathrel{+}= \text{"halluc"}$
    \EndIf
    \If{$q.\text{sol}\neq\varnothing \wedge q.\text{text}=\varnothing$}
      \State \textbf{delete} $q.\text{sol}$
    \EndIf
  \EndFor
\EndFor
\State $c \leftarrow \max(0,\;1-0.15|W|-\ldots)$;\; flag if $c<0.70$
\end{algorithmic}
\end{algorithm}
\end{minipage}

\subsubsection{Complexity and Correctness Analysis}

\begin{table}[H]
\caption{Complexity and correctness summary for the two core algorithms.}
\label{tab:complexity}
\centering\small
\begin{tabular}{lp{3.8cm}p{2.8cm}p{5cm}}
\toprule
\textbf{Algorithm} & \textbf{Time} & \textbf{Space} & \textbf{Correctness} \\
\midrule
HybridRetrieve
  & $O(n{\cdot}d)$ per query ($n$=chunks, $d$=768); $O(\sqrt{n}{\cdot}d)$ with IVF-PQ FAISS
  & $O(n{\cdot}d)$ index + $O(n)$ BM25 ($\approx$70~MB for $n{=}12{,}000$)
  & RRF with $k{=}60$ is Pareto-optimal over single rankers \cite{cormack2009rrf};
    diversity cap (${\leq}2$ chunks/source) prevents source monopoly. \\
\addlinespace
GroundingValidate
  & $O(N{\cdot}M)$ ($N$=questions, $M$=markdown length); ${\approx}75{,}000$ char comparisons in practice ($<$1~ms)
  & $O(N)$ for warning list
  & \textit{Soundness}: zero false negatives on orphan cleanup.
    \textit{Completeness}: probe may miss fragments inside table cells;
    asymmetric policy maximises recall. \\
\bottomrule
\end{tabular}
\end{table}

\subsubsection{Prompt Design}

A key engineering decision is the \textbf{separation of style and content} in the
generation prompt --- two distinct labeled blocks prevent topic blending
(Failure~F3 from the prototype).

\begin{figure}[H]
\noindent
\begin{minipage}[t]{0.48\textwidth}
\begin{lstlisting}[basicstyle=\ttfamily\tiny, breaklines=true,
  frame=single, backgroundcolor=\color{gray!6}]
[SYSTEM] Generate ONE MCQ at {bloom_level}
for subject: {subject}.

[STYLE ANCHORS]
Ex1: {question} | choices:[..] | correct:{i}
Ex2: {question} | choices:[..] | correct:{i}
(3-6 examples: same subject + Bloom level)

[COURSE MATERIAL - RAG]
{retrieved_chunks} <- HybridRetrieve()

[CONSTRAINTS]
- 4 choices (A-D); exactly 1 correct.
- Distractors: common misconceptions.
- Stem self-contained.
- Output JSON (enforced by instructor).

[OUTPUT SCHEMA]
{question_text, choices, correct_index,
 explanation, bloom_level, estimated_points}
\end{lstlisting}
\end{minipage}
\hfill
\begin{minipage}[t]{0.48\textwidth}
\begin{lstlisting}[basicstyle=\ttfamily\tiny, breaklines=true,
  frame=single, backgroundcolor=\color{gray!6}]
[SYSTEM] Evaluate the question strictly
against course material. No ext. knowledge.

[COURSE MATERIAL - fresh RAG]
{fresh_context} <- anti-circularity

[QUESTION TO EVALUATE]
{question_json}

[EVALUATION AXES]
1. factual_ok   : answer derivable?
2. level_correct: Bloom level matches?
3. quality      : good|needs_revision|reject
4. diff_fair    : point value appropriate?

[OUTPUT SCHEMA]
{factual_ok, level_correct,
 corrected_level, quality,
 diff_fair, estimated_minutes, comment}
\end{lstlisting}
\end{minipage}
\caption{Left: MCQ generation prompt structure. Right: post-generation
4-axis quality evaluation prompt. Both prompts are schema-enforced
via the \texttt{instructor} library.}
\label{lst:mcq_prompt}
\end{figure}
